\cleardoublepage\phantomsection\addcontentsline{toc}{chapter}{参考文献}
\begin{thebibliography}{48}
% 按全篇首次引用顺序排列。
\interlinepenalty=10000

\bibitem{appliedprecisioncvd}
Applied Materials. Producer XP Precision CVD.
n.d.
\url{https://www.appliedmaterials.com/us/en/product-library/producer-xp-precision-cvd.html} (accessed September 7, 2026).

\bibitem{pan2018survey}
Pan CR, Zhou MC, Qiao Y, Wu NQ. Scheduling cluster tools in semiconductor manufacturing: Recent advances and challenges.
IEEE Trans Autom Sci Eng 2018;15:586--601.
\url{https://doi.org/10.1109/TASE.2016.2642997}.

\bibitem{gu2024multifinger}
Gu L, Wu N, Qiao Y, Zhang S. Scheduling of {Multi-finger-robotic Cluster Tools} with {Multi-space Process Modules}.
In: 2024 International Conference on Networking, Sensing and Control ({ICNSC}). IEEE; 2024.
\url{https://doi.org/10.1109/ICNSC62968.2024.10760109}.

\bibitem{perkinson1994throughput}
Perkinson TL, McLarty PK, Gyurcsik RS, Cavin RK. Single-wafer cluster tool performance: An analysis of throughput.
IEEE Trans Semicond Manuf 1994;7:369--73.
\url{https://doi.org/10.1109/66.311340}.

\bibitem{rostami2002optimal}
Rostami S, Hamidzadeh B. Optimal scheduling techniques for cluster tools with process-module and transport-module residency constraints.
IEEE Trans Semicond Manuf 2002;15:341--9.
\url{https://doi.org/10.1109/TSM.2002.801379}.

\bibitem{wu2010closedform}
Wu NQ, Zhou MC. A closed-form solution for schedulability and optimal scheduling of dual-arm cluster tools with wafer residency time constraint based on steady schedule analysis.
IEEE Trans Autom Sci Eng 2010;7:303--15.
\url{https://doi.org/10.1109/TASE.2008.2008633}.

\bibitem{qiao2014revisiting}
Qiao Y, Wu N, Zhou MC. Scheduling of dual-arm cluster tools with wafer revisiting and residency time constraints.
IEEE Trans Ind Inform 2014;10:286--300.
\url{https://doi.org/10.1109/TII.2013.2272702}.

\bibitem{qiao2015thesis}
乔岩. 半导体制造中组合设备的建模、优化调度与控制 [D]. 广州：广东工业大学，2015.

\bibitem{zhou2022thesis}
周家忠. 基于时间Petri网离散事件系统的优化调度 [D]. 西安：西安电子科技大学，2022.

\bibitem{yuan2023survey}
袁凤连，黄波，王际鹏，潘春荣. 基于Petri网的组合设备建模与调度综述. 自动化学报，2023，49(5)：929--948.
\url{https://doi.org/10.16383/j.aas.c210951}.

\bibitem{wang2020multiple}
Wang J, Hu H, Pan C, Zhou Y, Li L. Scheduling dual-arm cluster tools with multiple wafer types and residency time constraints.
IEEE/CAA J Autom Sin 2020;7:776--89.
\url{https://doi.org/10.1109/JAS.2020.1003150}.

\bibitem{ko2021balance}
Ko SG, Yu TS, Lee TE. Wafer delay analysis and workload balancing of parallel chambers for dual-armed cluster tools with multiple wafer types.
IEEE Trans Autom Sci Eng 2021;18(3):1516--1526.
\url{https://doi.org/10.1109/TASE.2021.3061140}.

\bibitem{yang2025miptwowafers}
Yang F, Li C, Zhen L, Zhang C, Wu N. Scheduling residency time-constrained single-armed cluster tools with two-wafer types via mixed integer programming model.
Expert Syst Appl 2025;288:128209.
\url{https://doi.org/10.1016/j.eswa.2025.128209}.

\bibitem{li2023urgent}
李金成，潘春荣. 面向单臂组合设备中紧急订单的调度. 信息与控制，2023，52(5)：626--637.
\url{https://doi.org/10.13976/j.cnki.xk.2023.2303}.

\bibitem{luo2026mask}
罗亚波，巫琳，张峰. 考虑换模操作的双臂组合设备终止暂态调度. 工业工程，2026，29(1)：120--132.
\url{https://doi.org/10.3969/j.issn.1007-7375.240320}.

\bibitem{yang2025noncyclic}
Yang F, Xu X, Yang Z, Li C, Zhen L, Wu N. Non-cyclic scheduling of single-armed cluster tools for processing two wafer types.
IEEE Trans Autom Sci Eng 2025;22:14133--14146.
\url{https://doi.org/10.1109/TASE.2025.3554273}.

\bibitem{kim2023tandemphd}
Kim T. Scheduling and control of tandem cluster tools [D]. Daejeon: Korea Advanced Institute of Science and Technology, 2023.
\href{https://library.kaist.ac.kr/search/detail/view.do?bibCtrlNo=1046815&flag=dissertation}{KAIST Library, DIE~23014}.

\bibitem{pan2024twin}
潘春荣，周浩，熊文清，崔煜，罗继亮. 晶圆驻留时间约束下双组合设备协同调度. 信息与控制，2024，53(6)：804--816.
\url{https://doi.org/10.13976/j.cnki.xk.2024.3457}.

\bibitem{jee2021thesis}
Jee YJ. Scheduling of multi-slot cluster tools with wafer delay constraints [D]. Daejeon: Korea Advanced Institute of Science and Technology, 2021.
\href{https://library.kaist.ac.kr/search/detail/view.do?bibCtrlNo=956468&flag=dissertation}{KAIST Library, DIE~21004}.

\bibitem{zou2024twospace}
Zou C, Wang G, Gu L, Li J, Wu N, Qiao Y. Scheduling of time-constrained single-arm cluster tools with two-space process modules in semiconductor manufacturing.
IEEE Robot Autom Lett 2024;9(7):6352--6359.
\url{https://doi.org/10.1109/LRA.2024.3396112}.

\bibitem{huang2025efem}
Huang B, Wu N, Qiao Y, Guo W. Scheduling single-arm cluster tools with an equipment front-end module subject to wafer residency time constraints.
IEEE Trans Autom Sci Eng 2025;22:9544--9558.
\url{https://doi.org/10.1109/TASE.2024.3507854}.

\bibitem{yu2020delay}
Yu TS, Lee TE. Wafer delay analysis and control of dual-armed cluster tools with chamber cleaning operations.
Int J Prod Res 2020;58:434--47.
\url{https://doi.org/10.1080/00207543.2019.1593547}.

\bibitem{qiao2022binary}
Qiao Y, Lu Y, Li J, Zhang S, Wu N, Liu B. An efficient binary integer programming model for residency time-constrained cluster tools with chamber cleaning requirements.
IEEE Trans Autom Sci Eng 2022;19:1757--71.
\url{https://doi.org/10.1109/TASE.2021.3122576}.

\bibitem{zhen2023cnclean}
镇璐，张晓琴，阳罚军. 带驻留时间约束及清洗工艺的晶圆制造设备调度研究. 系统工程理论与实践，2023，43(8)：2395--2411.
\url{https://doi.org/10.12011/SETP2022-2338}.

\bibitem{li2022conditioncleaning}
Li C, Yang F, Zhen L. Efficient scheduling approaches to time-constrained single-armed cluster tools with condition-based chamber cleaning operations.
Int J Prod Res 2022;60:3555--68.
\url{https://doi.org/10.1080/00207543.2021.1926568}.

\bibitem{zhen2025cleaning}
Zhen L, Zhang B, Yang F, Zhou M. Efficient scheduling of residency-time-constrained dual-armed cluster tools with condition-based chamber cleaning operations.
IEEE Trans Autom Sci Eng 2025;22:5000--5011.
\url{https://doi.org/10.1109/TASE.2024.3414843}.

\bibitem{yang2025twoclean}
Yang F, Li C, Zhang C, Wu N. Optimal scheduling of single-armed cluster tools with two wafer types and chamber cleaning.
IEEE Trans Syst Man Cybern Syst 2025;55(8):5607--5618.
\url{https://doi.org/10.1109/TSMC.2025.3571733}.

\bibitem{chen2003dummy}
Chen HC, Lee CEC. Pull system for control and dummy wafers.
Int J Adv Manuf Technol 2003;22:805--818.
\url{https://doi.org/10.1007/s00170-002-1531-8}.

\bibitem{yang2009dummyphd}
楊懿淑. 需求不確定下控檔片之管理 [D]. 新竹：國立交通大學，2009.
\href{https://ir.lib.nycu.edu.tw/bitstream/11536/68667/1/380101.pdf}{学位论文全文}.

\bibitem{zhu2023regulation}
Zhu Q, Li B, Hou Y, Li H, Wu N. Scheduling dual-arm multi-cluster tools with regulation of post-processing time.
IEEE/CAA J Autom Sin 2023;10(8):1730--1742.
\url{https://doi.org/10.1109/JAS.2023.123189}.

\bibitem{kim2026flow}
Kim MC, Yu TS. Wafer flow time control in robotized semiconductor manufacturing tools.
IEEE Trans Syst Man Cybern Syst 2026;56(4):2333--2346.
\url{https://doi.org/10.1109/TSMC.2025.3649646}.

\bibitem{kim2022feedback}
Kim C, Yu TS, Lee TE. Feedback control of cluster tools: Stability against random time disruptions.
IEEE Trans Autom Sci Eng 2022;19(3):2008--2015.
\url{https://doi.org/10.1109/TASE.2021.3070346}.

\bibitem{pan2024variation}
Pan C, Lai Y, Qiao Y, Wu N, Luo X. Cycle time analysis for cluster tools with parallel process chambers, processing time variation, and chamber cleaning operations.
Expert Syst Appl 2024;255:124471.
\url{https://doi.org/10.1016/j.eswa.2024.124471}.

\bibitem{kim2025adaptivecluster}
Kim HJ, Lee JH. Scheduling cluster tools for concurrent processing: Deep reinforcement learning with adaptive search.
IEEE Trans Autom Sci Eng 2025;22:3783--96.
\url{https://doi.org/10.1109/TASE.2024.3399818}.

\bibitem{kim2025noncyclicdeadlock}
Kim D, Kim HJ. Noncyclic scheduling of cluster tools using deep reinforcement learning.
IEEE Trans Autom Sci Eng 2025;22:24009--23.
\url{https://doi.org/10.1109/TASE.2025.3632534}.

\bibitem{lai2026front}
Lai Y, Qiao Y, Chen L, He Y, Zhang S, Li J, et~al. A deep reinforcement learning-based method for non-cyclic scheduling of single-armed cluster tools with an equipment front-end module.
Adv Eng Inform 2026;74:104718.
\url{https://doi.org/10.1016/j.aei.2026.104718}.

\bibitem{huang2022cutranking}
Huang Z, Wang K, Liu F, Zhen HL, Zhang W, Yuan M, et~al. Learning to select cuts for efficient mixed-integer programming.
Pattern Recognit 2022;123:108353.
\url{https://doi.org/10.1016/j.patcog.2021.108353}.

\bibitem{paulus2022looking}
Paulus MB, Zarpellon G, Krause A, Charlin L, Maddison C. Learning to cut by looking ahead: Cutting plane selection via imitation learning.
In: Proceedings of the 39th International Conference on Machine Learning. PMLR; 2022, p. 17584--600.
\url{https://proceedings.mlr.press/v162/paulus22a.html}.

\bibitem{wang2023hem}
Wang Z, Li X, Wang J, Kuang Y, Yuan M, Zeng J, et~al. Learning cut selection for mixed-integer linear programming via hierarchical sequence model.
In: Proceedings of the Eleventh International Conference on Learning Representations. 2023.
\url{https://openreview.net/forum?id=Zob4P9bRNcK}.

\bibitem{wang2024hemset}
Wang J, Wang Z, Li X, Kuang Y, Shi Z, Zhu F, et~al. Learning to cut via hierarchical sequence/set model for efficient mixed-integer programming.
IEEE Trans Pattern Anal Mach Intell 2024;46(12):9697--9713.
\url{https://doi.org/10.1109/TPAMI.2024.3432716}.

\bibitem{zhou2016ccr}
周炳海，黎明. 基于CCR带并行腔双集束型设备调度方法. 北京航空航天大学学报，2016，42(7)：1361--1367.
\url{https://doi.org/10.13700/j.bh.1001-5965.2015.0538}.

\bibitem{fischetti2003local}
Fischetti M, Lodi A. Local branching.
Math Program 2003;98:23--47.
\url{https://doi.org/10.1007/s10107-003-0395-5}.

\bibitem{danna2005rins}
Danna E, Rothberg E, Le Pape C. Exploring relaxation induced neighborhoods to improve MIP solutions.
Math Program 2005;102:71--90.
\url{https://doi.org/10.1007/s10107-004-0518-7}.

\bibitem{hendel2022alns}
Hendel G. Adaptive large neighborhood search for mixed integer programming.
Math Program Comput 2022;14:185--221.
\url{https://doi.org/10.1007/s12532-021-00209-7}.

\bibitem{liu2022thesis}
Liu D. Machine learning algorithms for combinatorial optimization [D]. Montréal: Polytechnique Montréal, 2022.
\url{https://publications.polymtl.ca/10460/}.

\bibitem{krause2014submodular}
Krause A, Golovin D. Submodular function maximization.
In: Tractability: Practical Approaches to Hard Problems. Cambridge University Press; 2014.
\url{https://las.inf.ethz.ch/bibtexbrowser.php?bib=laspub.bib\&key=krause14survey}.

\bibitem{achterberg2009scip}
Achterberg T. {SCIP}: Solving constraint integer programs.
Math Program Comput 2009;1:1--41.
\url{https://doi.org/10.1007/s12532-008-0001-1}.

\bibitem{bestuzheva2023scip}
Bestuzheva K, Besan{\c{c}}on M, Chen WK, Chmiela A, Donkiewicz T, van Doornmalen J, et~al. Enabling research through the {SCIP} optimization suite 8.0.
ACM Trans Math Softw 2023;49:22:1--22:21.
\url{https://doi.org/10.1145/3585516}.

\end{thebibliography}
